% Generated draft: requires mathematical and source review.
\documentclass[11pt]{article}
\usepackage[margin=25mm]{geometry}
\usepackage{fontspec}
\setmainfont{Latin Modern Roman}
\usepackage{amsmath,amssymb,mathtools,unicode-math}
\setmathfont{Latin Modern Math}
\usepackage{enumitem,needspace,xurl,hyperref}
\setlength{\parindent}{0pt}
\setlength{\parskip}{5pt}
\setlength{\emergencystretch}{3em}
\newcommand{\N}{\mathbb{N}}
\newcommand{\Z}{\mathbb{Z}}
\newcommand{\Q}{\mathbb{Q}}
\newcommand{\R}{\mathbb{R}}
\newcommand{\C}{\mathbb{C}}
\newcommand{\F}{\mathbb{F}}
\newcommand{\A}{\mathbb{A}}
\newcommand{\PP}{\mathbb{P}}
\newcommand{\E}{\mathbb{E}}
\newcommand{\eps}{\varepsilon}
\newcommand{\dd}{\,\mathrm{d}}
\begin{document}
\section*{UnsolvedMath ID: 10400104}
\textbf{Source problem code:} catalogue entry
\subsection{Definitions and mathematical statement}
Let \(X\) be the source's quandle. A quandle operation is idempotent, has bijective right translations, and satisfies right self-distributivity. For each \(X\)-module \(A\) in the source's coefficient category, the task is to compute the second quandle cohomology group
\[
 H_Q^2(X;A).
\]
Once its coefficient cochain complex is specified, this means identifying the quotient of normalized 2-cocycles by 2-coboundaries,
\[
 H_Q^2(X;A)=\ker(\delta:C_Q^2(X;A)\to C_Q^3(X;A))\big/
 \operatorname{im}(\delta:C_Q^1(X;A)\to C_Q^2(X;A)),
\]
not merely giving the tautological kernel/image notation or isolated cocycles.

The excerpt does not define \(X\), the \(X\)-module structure on \(A\), the coefficient action or differential, and the normalization relations. Coefficients may carry nontrivial quandle-dependent structure, so ordinary trivial abelian coefficients cannot be substituted. Quandle cohomology must likewise be distinguished from rack cohomology before imposing degeneracy relations.

The quantifier is every admissible \(X\)-module \(A\), with no finiteness or field restriction supplied. A calculation only of rational dimensions or only for trivial modules would therefore address a narrower task. The surrounding source is needed to tell whether \(X\) is fixed earlier or ranges over a specified class, and to define the actual cochain complex.
\subsection{Short English statement}
Compute degree-two quandle cohomology for every module over the intended quandle, retaining its coefficient action and torsion. The source excerpt omits the quandle/module definitions and the resulting differential.
\subsection{Sources}
\begin{itemize}
\item Local UnsolvedMath registry, permanent ID 10400104, source code catalogue entry; no proposer or proposal year supplied; update 2026-08-06.
\item Ohtsuki, \emph{Problems on invariants of knots and 3-manifolds} (2002), Problem 5.8, PDF page 92: \url{https://msp.org/gtm/2002/04/gtm-2002-04-024s.pdf}. Registry access date 2026-07-29; PDF font-structure extraction. Provenance distinguishes 231 numbered bold headings. Current primary-literature status and redistribution terms require review.
\item The supplied summary repeats the module-coefficient calculation question without precise additional references or verified computations. Adjacent local records do not restore the omitted module/cochain definitions.
\end{itemize}
\end{document}
